% Source: sources/returns/S6_04_MAPS_MANIFOLDS/sections/manifolds.tex:8205–8473
\paragraph{Theorem D.2.41 (Proved: every positive weight by integral bar chains).}
The restriction on $\gcd(a,d)$ in D.2.37 and D.2.40 can
be removed. Every positive iterate $Q\mapsto dQ$ of
every one of the $172$ natural initial rows has an
exact integral homology algorithm, for all $495$
original/Mumford presentation patterns, every free
finite-good affine logarithm, and arbitrary regular
logarithms and clutching integer. Star caps remain
original. The full Seifert fundamental group and
the free ranks (D.WP8)--(D.WP9) are unchanged.

Here is the required all-parity correction to the
full matrix (D.OW2)--(D.OW4). Index the boundary
matrices as $A_0,\ldots,A_r$, with
$A_0\cdots A_r=I$, and use the first $r$ base loops
as independent generators. All $A_i$ below mean
the actual weighted matrices $A_i(d)$.
For $u,v\in\Lambda=\mathbb Z^4$ set
\[
 p(u,v)=\sum_{\alpha<\beta}u_\alpha v_\beta\,
                  e_\alpha\wedge e_\beta,\qquad
 D_A(u,v)=p(Au,Av)-(\Lambda^2A)p(u,v).
                                                        \tag{D.BC1}
\]
The bilinear map $D_A$ is symmetric. Define its
integral quadratic primitive by
\[
 f_A(u)=-\sum_\alpha D_A(e_\alpha,e_\alpha)
                       {u_\alpha\choose2}
        -\sum_{\alpha<\beta}D_A(e_\alpha,e_\beta)
                       u_\alpha u_\beta .
                                                        \tag{D.BC2}
\]
Negative coordinates use the same integral polynomial
$\binom n2=n(n-1)/2$.
Define the peripheral correction
\[
 \begin{split}
 c_i(u)&=f_{A_i}(u) \quad (i<r),\\
 c_r(u)&=-\sum_{j<r}
       f_{A_j}(A_{j+1}\cdots A_{r-1}A_r u).
 \end{split}                                            \tag{D.BC3}
\]
The transported reversed word in the second line
is part of the formula.

Keep every fibre relation (D.OW2) and every full
saturated cap kernel (D.OW4). Replace the five
finite or regular logarithmic normal columns by
\[
 \begin{gathered}
 \left(\sum_{\nu=0}^{m-1}c_i(A_i^\nu\lambda)
                  -\lambda\wedge v,\quad J_iN_m\lambda\right)
       \quad(\lambda=e_1,e_2,\delta,\gamma),\\
                       (c_i(v),J_iv).
 \end{gathered}                                          \tag{D.BC4}
\]
At every other cap the fixed-normal columns are
\[
                         (c_i(v),J_iv),
                         \qquad v\text{ in a basis of }\operatorname{Fix} A_i.
                                                        \tag{D.BC5}
\]
The finite freeness test is still exactly (D.WP5).
If $A_i$ has order $h\mid m$, the sum in (D.BC4)
is $(m/h)\sum_{\nu=0}^{h-1}c_i(A_i^\nu\lambda)$;
thus this is a finite algorithm for unbounded orders.
Let $T$ be the torsion cokernel of this full matrix.
Then the complete integral homology is (D.WP9)
with this $T$, for every $d>0$, including all
shared even factors.

\begin{proof}
\emph{An integral chain comparison with an explicit marking.}
Use the normalized bar chains of $\Lambda$ with
trivial coefficients. Their low-dimensional
boundaries are
\[
 \partial[u]=0,\qquad
 \partial[u|v]=[v]-[u+v]+[u].
                                                        \tag{D.BC6}
\]
The bar complex is a functorial chain model of
$B\Lambda$; see the standard resolution
\cite[Chapter I, Section 5]{manBrown82}.
We give all the comparison formulas used here.
The maps
\[
             p_1([u])=u,\qquad p_2([u|v])=p(u,v)
                                                        \tag{D.BC7}
\]
annihilate boundaries in the appropriate degree.
In degree two the cycle $[u|v]-[v|u]$ maps to
$u\wedge v$. The coordinate subtori give an
integral basis, so $p_2$ induces exactly
$H_2(T^4;\mathbb Z)\simeq\Lambda^2\Lambda$.

Antisymmetrizing (D.BC1) gives zero, proving
the symmetry of $D_A$. The binomial identity
$\binom{x+y}{2}=\binom{x}{2}+\binom{y}{2}+xy$
then gives the exact integral identity
\[
 f_A(u)+f_A(v)-f_A(u+v)
             =p(Au,Av)-(\Lambda^2A)p(u,v).
                                                        \tag{D.BC8}
\]
Equivalently, extending $f_A([u])=f_A(u)$
linearly on bar one-chains,
$f_A\partial=p_2A-(\Lambda^2A)p_2$ on bar
two-chains.

The regular piece retracts to its torus bundle
over a rose with $r$ petals. Its integral bar
mapping-cone complex has
\[
 {\cal C}_n=B_n(\Lambda)\oplus
                   \bigoplus_{i<r}B_{n-1}(\Lambda),
 \qquad
 \partial(c;g_i)=
 \left(\partial c+\sum_{i<r}(A_i-I)g_i;\
                                     -\partial g_i\right).
                                                        \tag{D.BC9}
\]
This is a model for the actual linear torus bundle,
not an assertion of equivariant formality. Its
universal cover is a torus-universal-cover bundle
over a tree, hence contractible. The bar model
has the same semidirect-product fundamental group,
and their maps on every peripheral torus agree
when their marked loops agree.

The sign in (D.BC9) uses the existing inverse
endpoint convention
$(x,1)\sim(A_i^{-1}x,0)$. Label a cylinder by its
terminal fibre chain $g$; its initial label is
$A_i g$. With fibre-first product orientation
its boundary is $(A_i-I)g$. This is the same cut
orientation used in (D.OW3).

In degree two define
\[
 P_2(c;g_i)=
 \left(p_2c+\sum_{i<r}f_{A_i}(g_i);\
                                  (p_1g_i)_{i<r}\right).
                                                        \tag{D.BC10}
\]
The target is the small complex whose degree-two
homology is
\[
 C\oplus K,\qquad
 C=\Lambda^2\Lambda/
          \sum_{i<r}(\Lambda^2A_i-I)\Lambda^2\Lambda,
 \quad
 K=\ker[A_0-I\ \cdots\ A_{r-1}-I].
                                                        \tag{D.BC11}
\]
No saturation or torsion quotient is taken in $C$.
Formula (D.BC8) proves that $P_2$ carries a
degree-three boundary to a small-complex boundary:
its fibre coordinate is
\[
 \sum_{i<r}\bigl(p_2(A_i-I)h_i-f_{A_i}\partial h_i\bigr)
       =\sum_{i<r}(\Lambda^2A_i-I)p_2h_i .
                                                        \tag{D.BC12}
\]
The contribution $p_2\partial c_3$ is zero.
Similarly $P_2$ commutes with the differential
to degree one. Thus it induces a homomorphism
on $H_2$.

Both cone homologies have the short exact
sequence with left term the degree-two fibre
coinvariant and right term $K$. On the left
$P_2$ is the integral isomorphism induced by
$p_2$, and on the right it is the identity
induced by $p_1$. The short five lemma proves
that (D.BC10) is an integral isomorphism
$H_2(V)\simeq C\oplus K$. This supplies the
missing two-primary marking directly. Neither
division by two nor descent of an invariant
differential form is used.

\emph{Fixed sweeps, norm sweeps, and the last boundary.}
If $A_i v=v$ and $i<r$, the literal fixed-loop
torus is the bar cycle with $g_i=[v]$ and
all other entries zero. Formula (D.BC10) sends
it to $(f_{A_i}(v),J_iv)$.
For an arbitrary loop $\lambda$ with
$A_i^m\lambda=\lambda$, the pushforward of
its torus in the $m$-fold mapping-torus cover
has edge chain
\[
                        \sum_{\nu=0}^{m-1}[A_i^\nu\lambda].
                                                        \tag{D.BC13}
\]
Its boundary telescopes to zero. Hence its
coordinates are the norm term in (D.BC4)
before subtracting $\lambda\wedge v$.

For the last boundary the base word is
$t_r=t_{r-1}^{-1}\cdots t_0^{-1}$.
The local bar one-chain $[u]$ maps to the
independent edge chains
\[
         -[A_{j+1}\cdots A_{r-1}A_ru],
                            \qquad 0\le j<r .
                                                        \tag{D.BC14}
\]
This identity is at the chain level: applying
$\sum(A_j-I)$ makes the bars telescope to
$[A_ru]-[u]$, with no leftover fibre two-chain.
For a fixed vector, its linear part is exactly
the previous marked $J_rv$. Applying
(D.BC10) gives the second line of (D.BC3).
Applying (D.BC14) to every summand of
(D.BC13) gives the full transported last-boundary
norm correction. In particular no reordering
of the last peripheral segments has been
silently discarded.

\emph{The actual cap maps.}
For a finite or regular logarithm the boundary
group has its normal meridian
$\mu=t_i^m v^{-1}$, with $A_i^m=I$ and
$A_i v=v$. The loops $\lambda,t_i^m,v$ lie
in one abelian subgroup. Bilinearity of its
torus classes therefore gives the exact
integral relation
\[
              [\lambda,\mu]=[\lambda,t_i^m]-[\lambda,v].
                                                        \tag{D.BC15}
\]
The first term is (D.BC13), and the second is
the fibre bivector $\lambda\wedge v$.
This proves the first four columns of (D.BC4)
even when their images have two-primary
torsion. The affine meridian generator gives
$[t_i,\mu]=[v,t_i]$, the fixed sweep already
computed, proving the fifth column. These
are exactly the full normal Thom generators
from D.2.29--30, including the torsion
normal-line cases.

Every fixed-normal torus at an unchanged cap
is still the same literal peripheral sweep,
so its image is (D.BC5). The fibre kernels
are exactly the saturated lattices proved
from the explicit local models in D.2.40.
The possible original multiplicative
correction $\delta\wedge u_2$ is zero in
the full global fibre coinvariant by the
$172$ unit elliptic certificates (D.OW7).
Thus the listed columns give the entire
marked kernel image at every allowed cap.

Finally the relative pair sequence (D.WP11)
has free final quotient for every weight;
its proof did not use the parity of $g$.
Its ranks depend only on the actual component
counts already proved in D.2.37. It follows
that the corrected full matrix computes
exactly the middle torsion, while the Seifert
group, free ranks and duality table stay
unchanged. All formulas are integral
polynomials, binomial polynomials, saturated
integer kernels and finite Smith operations.
They impose no upper bound on $d$, the
logarithmic orders or affine coordinates.
\end{proof}
